\documentclass{subfiles}

\begin{document}
    Berdasarkan hasil praktikum pada Modul {\arabic{chapter}} tentang {\leftmark}, dapat diambil kesimpulan bahwa:
    \begin{enumerate}
        \item {
            \textit{Array} merupakan tipe data yang ada diberbagai bahasa pemrograman. Biasanya tipe data \textit{array} terdiri dari kumpulan-kumpulan nilai-nilai yang bertipe data sama. Sebagai contoh, sebuah vektor dapat disebut sebagai sebuah \textit{array} karena berupa kumpulan nilai-nilai yang bertipe data sama.
        }
        \item {
            \textit{Array} dalam bahasa pemrograman C++ dapat diinisialisasikan dengan menggunakan \textit{syntax} tipe data diikuti nama nama variabel yang dinginginkan yang diikuti oleh tanda kurung kotak yang di dalamnya terdapat ukuran dari {\sarray} yang diinginkan. Jumlah data pada \textit{array} dapat dikosongkan dan {\compiler} akan menebak jumlah data secara otomatis.
        }
        \item {
            \textit{String} merupakan tipe data yang tersusun dari kumpulan karakter. Terdapat berbagai fungsi-fungsi yang bisa digunakan untuk memanipulasi \textit{string} dalam C++ seperti \textit{size}, \textit{length}, \textit{clear} dan lain-lain.
        }
    \end{enumerate}
\end{document}
